% !TEX root = 00_main_aaai_2027.tex
\begin{abstract}
% CS197 move 1: problem motivation (failure mode) -- association does not identify probability forecasts.
A judge can predict the correct mean rating while assigning the wrong conditional distribution to every input.
% CS197 move 2: assumption (setup bit) -- scope is correlation alone, not all validation methods.
Correlation between scalar judge and human scores cannot distinguish such forecasts.
% CS197 move 3: contribution (bit flip) -- identify the desired target explicitly.
We argue for evaluating the conditional rating distribution with proper probability losses, alongside scalar association and agreement diagnostics.
% CS197 move 4: instantiation (bit flip instantiation) -- controlled synthetic identification, not natural-human validation.
A binomial construction gives five candidate distributions with identical conditional means and distinct expected cross-entropy, Brier score and ranked probability score.
% CS197 move 5a: evaluation -- provenance: experiment02/results/summary.json, m=1 ce.
With one synthetic rating per item, cross-entropy selects the true distribution in 249 of 256 replications.
% CS197 move 5b: diagnostic -- provenance: experiment02/results/selection_rows.csv, all scalar criteria.
Pearson, Spearman, Kendall and scalar mean squared error retain all five candidates as ties throughout.
% CS197 move 5c: limitations -- include every tie, preserve contradictory historical evidence.
We enumerate these ties, reproduce the older marginal-histogram example, and retain inconclusive historical theorem-judge evidence.
% CS197 move 6: implication (north star) -- narrow supported thesis.
The result supports distribution-aware evaluation without implying human generalization, universal finite-sample superiority, or that psychometric methods are unnecessary.
\end{abstract}
